Byte-pair encoding (BPE) tokenizers are trained once and then frozen into every model built on them, yet the order of their merges is almost never varied.
We ask whether merge order matters, both the order in which merges are \emph{learned} and the priority with which they are \emph{applied}.
We train six byte-level BPE tokenizers with the same 8,192-entry vocabulary size on the same corpus; they differ only in merge order.
Two of them keep the exact vocabulary of standard frequency BPE and change only merge priority.
For each tokenizer we pre-train identical 23M-parameter language models on the same 200\,MB of general and biomedical text, and then fine-tune them on three classification tasks.
The answer depends on what we measure.
For language modeling, merge order clearly matters (all 10 comparisons with Holm-adjusted $p \le 0.032$): the four non-frequency orders raise bits-per-byte by 1.8--5.8\% on general text and by 2.8--9.8\% on biomedical text.
Even reordering the \emph{same} merges costs 1.8--5.7\%, because short merges fire first and block long ones, leaving 59--79\% of the vocabulary unused.
Compression explains this effect: bytes-per-token perfectly rank-predicts general-text bits-per-byte across tokenizers (Spearman $\rho=-1.0$).
For fine-tuned classification, merge order shows no detectable effect: all tokenizers are within $-0.8\%$ to $+1.6\%$ of standard BPE, and none of the 15 comparisons is significant.
A domain-weighted merge order gives a small domain gain (0.8\% bits-per-byte) and no classification gain.
Merge order is therefore not a cosmetic detail for generative modeling, but small fine-tuned classifiers are robust to it.
